% Generated by tools/edn_latex.py from reports/edn.md.
% Do not edit: change reports/edn.md and run `make edn`.
\ednentry{
  id = {EDN-17},
  title = {Deployment target is the FIB Virtech VM},
  date = {2026-09-29},
  milestone = {M4: Model Deployment},
  topic = {ML System Design, Deployment},
  participants = {@lukas2510},
  decision = {The component is deployed on the FIB Virtech VM the course provides (4 GB RAM, 20 GB disk, CPU-only, one semester). The project brief no longer lists the target as an open point, and the requirements say so once instead of assuming it in every target.},
  rationale = {\emph{Alternatives considered:}
\begin{itemize}[leftmargin=*, topsep=2pt]
\item \textbf{Option A (chosen): the FIB Virtech VM.}
\newline Pros: free and provided for the course, so no payment details and no personal cloud account for a semester project; 4 GB is enough with the budgets NFR-04 already sets (drift job scheduled in its own container, Prometheus capped by retention, logs capped by rotation, 4 GB disk kept free); every target in the requirements is written for this machine.
\newline Cons: tight on RAM, a single unattended host with no redundancy, and nobody on the team has access yet.
\item \textbf{Option B: another cloud (AWS, DigitalOcean, Oracle Always Free or similar).}
\newline Pros: more headroom, and a free ARM tier exists at one provider.
\newline Cons: either paid or tied to a private account with payment details; the grading rewards nothing for the hosting choice; NFR-02 to NFR-04 would have to be re-derived for different hardware, and an ARM tier would also mean multi-architecture images.
\end{itemize}
\par
\emph{Why:} The requirements already named this VM in eight places (NFR-02, NFR-04, NFR-05, NFR-09, NFR-12, NFR-13 and the network warning), so the decision had been made by writing rather than by deciding, while the brief still listed the target as open. The reason the brief gave for keeping it open, that the stack is tight on 4 GB, has since been planned for in NFR-04, so the open point was stale. What stays genuinely open is not the target but what the VM's network allows, which the requirements already track with a fallback per requirement. Access is a task, not a decision, and it has a date: if nobody has access at the M4a lab, we raise it with the teachers.},
  ai = {Information seeking, Alternative assessment, Recommendation},
  response = {Accepted},
  assessment = {While reviewing PR \#19, AI reported the brief and the requirements as contradicting each other on the deployment target. Asked to look again, it corrected its own framing: the \texttt{[proposed]} marker in the requirements refers to the performance numbers, not to the machine, and the requirements treat the VM as settled throughout, so the real issue was an implicit decision rather than a disagreement between two pages. It also pointed out that the brief's stated reason for leaving the point open is already covered by NFR-04. Lukas kept the change small and deliberately left the performance targets \texttt{[proposed]} until the first M4 load test.},
  aievidence = {Claude Code session on 2026-09-29: after Lukas asked (translated from German) ``analyse that again, then explain it to me and give me a reasoned recommendation on what the best way is'' about the deployment target, AI recommended fixing the VM as the target with a dated trigger for the missing access, and Lukas replied ``yes do that''.},
  evidence = {\href{https://github.com/mlops-2627q1-mds-upc/MLOPS_Recommenditos/blob/main/docs/docs/specification.md}{Specification} section 2 and NFR-04; \href{https://github.com/mlops-2627q1-mds-upc/MLOPS_Recommenditos/blob/main/docs/docs/project-brief.md}{project brief} section 5; \href{https://github.com/mlops-2627q1-mds-upc/MLOPS_Recommenditos/pull/19}{PR \#19}.},
}
